\section{Предобработка}

Используется стандартный для CIFAR-10 набор преобразований:
\begin{itemize}
    \item train: \texttt{RandomCrop(32, padding=4)}, \texttt{RandomHorizontalFlip()}, \texttt{Normalize(mean,std)};
    \item val/test: \texttt{ToTensor()} + \texttt{Normalize(mean,std)}.
\end{itemize}

\subsection{Почему именно так}
\begin{itemize}
    \item Аугментации на train уменьшают переобучение и повышают устойчивость оценки параметров $\theta$.
    \item Единая нормализация train/val нужна для сопоставимости лосса в surface-оценке.
\end{itemize}

\subsection{Риски leakage}
В проекте нет data leakage между train и val: split берётся из официального CIFAR-10 протокола torchvision, а валидационный лосс для surface считается на test/val части без train-аугментаций.
